% Generated with genpdf
% Domain Interno
% Title Glossario
\documentclass{scalatekids-article}
\begin{document}
\lfoot{Glossario 0.0.3}
\newgeometry{top=3.5cm}
\begin{titlepage}
  \begin{center}
    \begin{center}
      \includegraphics[width=10cm]{sklogo.png}
    \end{center}
    \vspace{1cm}
    \begin{Huge}
      \begin{center}
        \textbf{Glossario}
      \end{center}
    \end{Huge}
    \vspace{11pt}
    \bgroup
    \def\arraystretch{1.3}
    \begin{tabular}{r|l}
      \multicolumn{2}{c}{\textbf{Informazioni sul documento}} \\
      \hline
      \setbox0=\hbox{0.0.1\unskip}\ifdim\wd0=0pt
      \\
      \else
      \textbf{Versione} & 0.0.4\\
      \fi
      \textbf{Redazione} & \multiLineCell[t]{Redattore}\\
      \textbf{Verifica} & \multiLineCell[t]{Verificatore}\\
      \textbf{Approvazione} & \multiLineCell[t]{Approvatore}\\
      \textbf{Uso} & Esterno\\
      \textbf{Lista di Distribuzione} & \multiLineCell[t]{ScalateKids\\Prof. Tullio Vardanega\\Prof. Riccardo Cardin}\\
    \end{tabular}
    \egroup
    \vspace{22pt}
  \end{center}
\end{titlepage}
\restoregeometry
\clearpage
\setcounter{page}{1}
\begin{flushleft}
  \vspace{0cm}
         {\large\bfseries Diario delle modifiche \par}
\end{flushleft}
\vspace{0cm}
\begin{center}
  \begin{tabular}{|l | l | l | l | l |}
    \hline
    Versione & Autore & Ruolo & Data & Descrizione \\
    \hline
    0.0.3 & Andrea Giacomo Baldan & & 2015-12-23 & Prima verifica\\
    \hline
    0.0.2 & Michael Munaro & & 2015-12-22 & Prima stesura voci riscontrate\\
    \hline
    0.0.1 & Marco Boseggia & & 2015-12-22 & Creazione scheletro del documento\\
    \hline
  \end{tabular}
\end{center}
\subsection{Scopo del documento}
Questo documento ha il compito di raccogliere tutti i termini che possono risultare sconosciuti o di difficile comprensione ad un lettore esterno oppure, semplicemente, concetti specifici che necessitano di essere esposti per poter meglio comprendere quanto viene detto nei documenti.
\subsection{Riferimenti} //da fare
\newpage
\section{Definizioni}

\glossaryLetter{A}

\glossDef{Actormodel} Modello matematico per la gestione della concorrenza computazionale proposto nel 1973 da Carl Hewitt.
All'interno del modello, gli "attori" vengono trattati come primitive universali e possono scmabiarsi messaggi tra di loro; in risposta ad un messaggio ricevuto, un attore può prendere decisioni, creare nuovi attori, inviare altri messaggi e determinare come rispondere al messaggio successivo.

\glossDef{Akka} E' uno strumento per la semplificazione e costruzione di applicazioni concorrenti e distribuite sulla JVM (Java Virtual Machine). Akka supporta diversi modelli per la concorrenza prediligendo, però, l'uso del modello ad attori.

\glossDef{Amazon DynamoDB} Database proprietario di tipo NoSQL e completamente gestito da Amazon.com all'interno del suo portfolio 'Amazon Web Services'.

\glossDef{Attività} Aggregazione di task (vedi sotto). %TODO

\glossDef{Attori}  vedi \textit{Actormodel}

\glossaryLetter{B}

\glossDef{Beacon} il beacon è un piccolo dispositivo BLE (acronimo per Bluetooth Low Energy, cioè bluethooth a basso consumo energetico) che trasmette piccole quantità di dati ogni secondo o quasi.  A livello di hardware, i beacon sono dispositivi BLE-compatible che trasmettono dati sotto il protocollo “beacon”.
A livello software, i beacon sono messaggi inviati da questi dispositivi trasmittenti, che sono poi rilevati e processati da un dispositivo ricevente, come un’app mobile. Nel caso degli iBeacon vengono trasmessi i valori Major e Minor che insieme all’UUID (acronimo di Universally Unique IDentifier), permettono l’identificazione di un iBeacon da parte dell’applicazione.(riferirsi a https://it.wikipedia.org/wiki/IBeacon e http://www.beaconitaly.it/cose-ibeacon/);

\glossDef{Brainstorming} Metodo decisionale in cui la ricerca della soluzione di un problema è effettuata mediante sedute intensive di dibattito e confronto delle idee e delle proposte espresse dai partecipanti.

\glossDef{Bundled, Software} Si tratta di software distribuito con un altro prodotto come parte di un hardware o come parte di un gruppo di pacchetti software venduti assieme.

\glossDef{Bug} in informatica, identifica un errore nella scrittura di un programma software, impedendone il corretto funzionamento e provocandone un comportamento imprevisto o comunque diverso da quello specificato dal produttore. Meno comunemente, il termine bug può indicare un difetto di progettazione in un componente hardware.
L'origine del nome (che in inglese significa "insetto") deriva da un malfunzionamento del computer Mark II nel 1947 dovuto ad una falena incastratasi tra i circuiti.

\glossaryLetter{C}

\glossDef{CamelCase (Notazione a Cammello)} è la pratica nata durante gli anni settanta di scrivere parole composte o frasi unendo tutte le parole tra loro, ma lasciando le loro iniziali maiuscole. La prima lettera può essere sia maiuscola che minuscola.
Esistono diverse varianti della CamelCase tra le quali ricordiamo la \textbf{UpperCamelCase}, ovvero una notazione a cammello in cui la prima lettera è obbligatoriamente maiuscola, e la \textbf{lowerCamelCase} dove la prima lettera è sempre minuscola.

\glossDef{CLI (Command Line Interface)} detta semplicemente \textit{riga di comando}, indica una tipologia di interfaccia utente caratterizzata da un'interazione di tipo testuale tra utente ed elaboratore.

\glossDef{Codice} Testo di un algoritmo di un programma scritto in un linguaggio di programmazione da un programmatore.

\glossDef{Commit} Nel contesto di computer science e gestione di dati, commit si riferisce all'idea di rendere permanenti i cambiamenti effettuati nel repository locale

\glossaryLetter{D}

\glossDef{Diagramma di Gantt} E' uno strumento di supporto alla gestione dei progetti.
Si tratta, sostanzialmente, di un diagramma costruito a partire da un asse orizzonatale, a rappresentazione delll'arco temporale totale (o parziale) del progetto, e da un asse verticale che rappresenta le mansioni o attività che costituiscono il progetto.
Lo scopo principale del diagramma di gantt è di dislocare temporalmente le attività per poterne rappresentare la durata, la sequenzialità e il parallelismo e poter confrontare le stime con i progressi.

\glossDef{Database} In informatica, il termine database (in italiano \textit{base di dati}) indica un insieme organizzato di dati.
Le informazioni contenute in un database sono strutturate e collegate tra loro secondo un particolare modello logico scelto dal progettista del database (ad esempio relazionale, gerarchico, key-value, ...).
Gli utenti si interfacciano con i database attraverso i cosiddetti \textit{query language} e grazie a particolari applicazioni software dedicate.

\glossDef{Design Pattern} In ingegneria del software, un design pattern descrive una soluzione generale a un problema di progettazione ricorrente, gli attribuisce un nome, astrae e identifica gli aspetti principali della struttura utilizzata per la soluzione del problema, identifica le classi e le istanze partecipanti e descrive quando e come può essere applicato. In breve definisce un problema, i contesti tipici in cui si trova e la soluzione ottimale allo stato dell'arte.

\glossDef{Dizionario (o Array Associativo)} Struttura dati che contiene una collezione di oggetti (o records).
Questi oggetti contengono diversi campi, ognuno dei quali contiene dei dati.
E' possibile accedere, gestire e salvare i diversi records tramite una chiave che li identifica in maniera univoca.

\glossDef{Docker} è un progetto open-source che automatizza il deployment di applicazioni all'interno di container software, fornendo un'astrazione addizionale grazie alla virtualizzazione a livello di sistema operativo di Linux.

\glossDef{Domain Specific Language (DSL)} E' un linguaggio di programmazione o un linguaggio di specifica dedicato a particolari problemi di un dominio, a una particolare tecnica di rappresentazione e/o a una particolare soluzione tecnica.

\glossaryLetter{E}

\glossaryLetter{F}

\glossDef{File} In informatica, viene utilizzato per riferirsi a un contenitore di informazioni/dati in formato digitale, tipicamente presenti su un supporto digitale di memorizzazione opportunamente formattato in un determinato file system.
Le informazioni scritte/codificate al suo interno sono leggibili solo tramite uno specifico software in grado di effettuare l'operazione.

\glossDef{Forum (internet)} Il termine "forum", nell'ambito dell'informatica, è utilizzato per indicare l'insieme delle sezioni di discussione in una piattaforma informatica dove gli utenti intrattengono conversazioni postando messaggi.
Un forum può essere di tipo gerarchico e con struttura ad albero in quanto possono esservi al suo interno diversi subforum, ognuno dei quali con un proprio topic (argomento di discussione).

\glossaryLetter{G}

\glossDef{Git} Git è un sistema software di controllo di versione distribuito, creato da Linus Torvalds nel 2005.
Esso permette, dunque, di tenere traccia delle modifiche e delle versioni apportate al codice sorgente di un software, senza la necessità di dover utilizzare un server centrale.
Con questo sistema gli sviluppatori possono collaborare individualmente e parallelamente offline su di un proprio ramo (\textit{branch}) di sviluppo, registrare le proprie modifiche (\textit{commit}) ed in seguito condividerle o unirle (\textit{merge}) a quelle di altri, il tutto, come detto in precedenza, senza bisogno del supporto di un server.

\glossDef{GitHub} Servizio Web di hosting per lo sviluppo di progetti software, he usa il sistema di controllo  versione Git.

\glossaryLetter{H}

\glossaryLetter{I}

\glossDef{IDE (Integrated development environment)} Un \textit{ambiente di sviluppo integrato} è un software che, in fase di programmazione, aiuta i programmatori nello sviluppo del codice sorgente di un programma. Spesso l'IDE aiuta lo sviluppatore segnalando errori di sintassi del codice direttamente in fase di scrittura, oltre a tutta una serie di strumenti e funzionalità di supporto alla fase di sviluppo e debugging.

\glossDef{IEC (International Electrotechnical Commission)} E' un'organizzazione internazionale per la definizione di standard in materia di elettricità, elettronica e tecnologie correlate, fondata a Londra nel 1906.

\glossDef{Instant messaging} E' una categoria di sistemi di comunicazione in tempo reale in rete, tipicamente Internet o una rete locale, che permette ai suoi utilizzatori lo scambio di brevi messaggi.

\glossDef{Integrazione Continua} Nell'ingegneria del software, l'integrazione continua (continuous integration in inglese, spesso abbreviato in CI) è una pratica che si applica in contesti in cui lo sviluppo del software avviene attraverso un sistema di versioning. Consiste nell'allineamento frequente dagli ambienti di lavoro degli sviluppatori verso l'ambiente condiviso (mainline).

\glossDef{IoT (Internet of Things)} è un neologismo riferito all'estensione di Internet al mondo degli oggetti e dei luoghi concreti.
L' \textit{Internet delle cose} è vista come una possibile evoluzione dell'uso della Rete: gli oggetti si rendono riconoscibili e acquisiscono intelligenza grazie al fatto di poter comunicare dati su se stessi e accedere ad informazioni aggregate da parte di altri[

\glossDef{ISO (International Organization for Standardization)} E' la più importante organizzazione a livello mondiale per la definizione di norme tecniche fondata nel 1947 a Ginevra, in Svizzera.
  L'ISO coopera strettamente con l'IEC (vedi sopra), responsabile per la standardizzazione degli equipaggiamenti elettrici.
  Le norme ISO sono numerate e hanno un formato del tipo ISO nnnn:yyyy - titolo, dove nnnn è il numero della norma, yyyy l'anno di pubblicazione e titolo è una breve descrizione della norma.

 \glossaryLetter{J}

  \glossaryLetter{K}

  \glossDef{Key-value Database} E' un paradigma di archivazione di dati progettato per archiviare, recuperare e gestire array associativi (vedi "Dizionario").

  \glossaryLetter{L}

  \glossDef{lowerCamelCase} vedi \textit{CamelCase}.

  \glossaryLetter{M}
  
  \glossDef{Macro} In informatica indica una procedura, ovvero un insieme di comandi o istruzioni, tipicamente ricorrente durante l'esecuzione di un programma. 

  \glossDef{Mailing List} E' un servizio/strumento offribile da una rete di computer verso vari utenti e costituito da un sistema organizzato per la partecipazione di più persone ad una discussione asincrona o per la distribuzione di informazioni utili agli interessati/iscritti attraverso l'invio di email ad una lista di indirizzi di posta elettronica di utenti iscritti.

  \glossDef{Modello ad attori} vedi "Actormodel".

  \glossDef{Modello relazionale} modello logico di rappresentazione o strutturazione dei dati di un database implementato su sistemi di 'gestione di basi di dati' (DBMS).
  Si basa sulla teoria degli insiemi e sulla logica del primo ordine ed è strutturato intorno al concetto matematico di relazione (detta anche tabella). Per il suo trattamento ci si avvale di strumenti quali il calcolo relazionale e l'algebra relazionale.

  \glossaryLetter{N}

  \glossDef{NoSQL} E' un movimento che promuove sistemi software dove la persistenza dei dati è caratterizzata dal fatto di non utilizzare il modello relazionale (vedi sopra), solitamente utilizzato dai database tradizionali.
  L'espressione NoSQL fa riferimento al linguaggio SQL, il più comune tra i linguaggi per l'interrogazione dei database relazionali, qui preso a simbolo dell'intero paradigma relazionale. Molto spesso nei linguaggi NoSQL viene sacrificato in favore dele prestazioni il paradigma ACID (atomicità, consistenza, isolamento, durabilità) in quanto, per il teorema CAP, noto anche come teorema di Brewer, è impossibile per un sistema informatico distribuito, fornire simultaneamente tutte e tre le seguenti garanzie:
  Coerenza (tutti i nodi vedono gli stessi dati nello stesso momento)
  Disponibilità (la garanzia che ogni richiesta riceva una risposta su ciò che sia riuscito o fallito)
  Tolleranza di partizione (il sistema continua a funzionare nonostante arbitrarie perdite di messaggi)
  Secondo il teorema, un sistema distribuito è in grado di soddisfare al massimo due di queste garanzie allo stesso tempo, ma non tutte e tre.

  \glossaryLetter{O}

  \glossDef{Organigramma} In un'organizzazione aziendale, con il termine organigramma si intende una formalizzazione di quali sono gli organi e le loro relazioni gerarchiche o funzionali all'interno di un smbiente di lavoro.

  \glossaryLetter{P}

  \glossDef{Plugin} Il plugin in campo informatico è un programma non autonomo che interagisce con un altro programma per ampliarne o estenderne le funzionalità originarie.
  A seconda dei programmi e delle piattaforme software, i plugin possono esser chiamati con sinonimi diversi: add-in, add-on e estensione.

  \glossDef{Prototipo} Il modello realizzato nell’ultima fase della progettazione e sperimentazione, e destinato a divenire il punto di partenza della produzione in serie
  
  \glossDef{Pull} %TODO
  
  \glossDef{Push} %TODO

  \glossaryLetter{Q}
  
  \glossDef{Query} termine che viene usato in informatica per indicare l'interrogazione di un database da parte di un utente.

  \glossaryLetter{R}

  \glossDef{Refactoring} In ingegneria del software, il refactoring (o code refactoring) è una tecnica per modificare la struttura interna di porzioni di codice senza modificarne il comportamento esterno applicata per migliorare alcune caratteristiche non funzionali del software.
  I vantaggi che il refactoring persegue riguardano in genere un miglioramento della leggibilità della manutenibilità, della riusabilità e dell'estendibilità del codice e la riduzione della sua complessità, eventualmente attraverso l'introduzione a posteriori di design pattern.

  \glossDef{Repository (Software Repository)} E' uno spazio di archiviazione da dove è possibile recuperare e scaricare dei pacchetti software.

  \glossaryLetter{S}

  \glossDef{Scala} E' un linguaggio di programmazione di tipo general-purpose studiato per integrare le caratteristiche e funzionalità dei linguaggi orientati agli oggetti e dei linguaggi funzionali. La compilazione di codice sorgente Scala produce Java bytecode per l'esecuzione su una JVM.

  \glossDef{Scaladoc} E' un applicativo incluso nel Scala Development Kit, utilizzato per la generazione automatica della documentazione del codice sorgente scritto in linguaggio Scala.

  \glossDef{Script} Il termine script, in informatica, designa un tipo particolare di programma, scritto in una particolare classe di linguaggi di programmazione, detti linguaggi di scripting.
  La distinzione tra un programma normale ed uno script non è netta ma solitamente gli script si distinguono per la loro complessità relativamente bassa, la mancanza di una propria interfaccia grafica e il fatto di essere utilizzati per svolgere mansioni accessorie e molto specifiche.
  
  \glossDef{Shell} in informatica, è la parte di un sistema operativo che permette agli utenti di interagire con il sistema stesso, impartendo comandi e richiedendo l'avvio di altri programmi.
  La shell può essere di tipo \textit{testuale} (vedi \textit{terminale}) o \textit{grafico}.

  \glossDef{SVN} Subversion (noto anche come svn, che è il nome del suo client a riga di comando) è un sistema di controllo versione per software distribuito come free software sotto licenza Apache.
  Gli sviluppatori software utilizzano Subversion per mantenere le versioni correnti e storiche di file quali codice sorgente, pagine web e documentazione.

  \glossaryLetter{T}
  
  \glossDef{Task} Compito assegnabile ad un singolo componente del gruppo di lavoro che dovrebbe essere ultimato entro un periodo di tempo prefissato o entro una deadline. %TODO

  \glossDef{Template} Trattasi di un modello che ha il compito di fornire la struttura di un oggetto (documento, sito web, ...) in cui si trovano spazi temporaneamente "bianchi" da riempire successivamente.
  
  \glossDef{Terminale} \textit{Shell} di tipo testuale che permette all'utente di interagire con il sistema operativo tramite un'interfaccia a riga di commando (vedi \textit{CLI})

  \glossDef{Topic} Tema di discussione, argomento.
  
  \glossDef{Travis CI} Servizio distribuito per l'integrazione continua usato per testare i progetti hostati su GitHub.

  \glossaryLetter{U}
  
  \glossDef{UI (User Interface)} L'interfaccia utente è ciò che si frappone tra una macchina e un utente, consentendo l'interazione tra i due. In ambito informatico sono tipicamente utilizzati i seguenti tipi di interfaccia utente: 
  \begin{itemize}
  	\item Interfaccia a riga di commando (CLI)
  	\item Interfaccia grafica (GUI)
  \end{itemize}

  \glossDef{UpperCamelCase} vedi \textit{CamelCase}.

  \glossDef{Unicode} sistema di codifica che assegna un numero univoco ad ogni carattere usato per la scrittura di testi, in maniera indipendente dalla lingua, dalla piattaforma informatica e dal programma utilizzato

  \glossDef{UTF-8 (Unicode Transformation Format, 8 bit)} E' una codifica dei caratteri Unicode (vedi sopra) in sequenze di lunghezza variabile di byte, creata da Rob Pike e Ken Thompson. UTF-8 usa gruppi di byte per rappresentare i caratteri Unicode, ed è particolarmente utile per il trasferimento tramite sistemi di posta elettronica a 8-bit.
  In particolare, UTF-8 usa da 1 a 4 byte per rappresentare un carattere Unicode.

  \glossaryLetter{V}

  \glossDef{Versionamento} Il versionamento (versioning), in informatica, è la gestione di versioni multiple di un insieme di informazioni.
  Gli strumenti software per il controllo versione sono ritenuti molto spesso necessari per la maggior parte dei progetti di sviluppo software.

  \glossDef{Virtual Machine} Il termine macchina virtuale (VM) indica un software che, attraverso un processo di virtualizzazione, crea un ambiente virtuale che emula tipicamente il comportamento di una macchina fisica grazie all'assegnazione di risorse hardware (porzioni di disco rigido, RAM e risorse di processamento) ed in cui alcune applicazioni possono essere eseguite come se interagissero con tale macchina.
  Tra i vantaggi vi è il fatto di poter offrire contemporaneamente ed efficientemente a più utenti diversi ambienti operativi separati, ciascuno attivabile su effettiva richiesta, senza sporcare il sistema fisico reale con il partizionamento del disco rigido oppure fornire ambienti clusterizzati (connessi tra loro) su sistemi server.

  \glossDef{VoIP (Voice over IP)} tecnologia che rende possibile effettuare una conversazione telefonica sfruttando una connessione Internet o una qualsiasi altra rete dedicata a commutazione di pacchetto che utilizzi il protocollo IP senza connessione per il trasporto dati.

  \glossaryLetter{W}

  \glossDef{Whiteboarding} E' una tecnica di condivisione di un notebook o di una whiteboard (Lavagna) virtuali nei quali più persone possono lavorare contemporaneamente per scambiarsi idee, bozze di grafici e schemi.
  Tutto questo viene fatto online così da permettere ai singoli utenti di vedere in tempo reale le modifiche fatte al documento.
  
  \glossDef{Wiki} Il wiki è un sito web, costruito appoggiandosi su una piattaforma o software collaborativo detto software wiki o semplicemente Wiki, che permette ai propri utenti di aggiungere, modificare o cancellare contenuti apertamente attraverso un browser web, in genere utilizzando un linguaggio di markup semplificato o un editor di testo online.
  Si tratta in altre parole di una raccolta di documenti ipertestuali che viene aggiornata dai suoi stessi utilizzatori e i cui contenuti sono sviluppati in collaborazione da tutti coloro che vi hanno accesso (contenuto generato dagli utenti), memorizzati normalmente su un database o un repository.

  \glossaryLetter{X}

  \glossaryLetter{Y}

  \glossaryLetter{Z}

\end{document}
